%! Introducing Vectors — Unit 1, Lesson 1.0 — Lesson Plan
\documentclass[10pt]{article}
\usepackage{linalg-article}
\usepackage{linalg-boxes}
\usepackage{graphicx}   % -article does NOT load graphicx; needed for thumbnails/screenshots
\graphicspath{{images/}}
\newcommand{\TallMath}[1]{$\displaystyle #1\rule[-1.4em]{0pt}{3.2em}$}
\newcommand{\vv}{\mathbf{v}}
\newcommand{\ww}{\mathbf{w}}
\newcommand{\UnitNumberName}{Unit 1: Vectors and Matrices \quad}
\newcommand{\LessonNumberName}{Lesson 1.0: Introducing Vectors --- Vocabulary and Length}

\begin{document}
\begin{center}
    {\Large\bfseries \CourseName: \SchoolYear \\
    {\small\bfseries \UnitNumberName \LessonNumberName}}
\end{center}

\begin{tcolorbox}[colback=blush, colframe=burgundy, boxrule=0.9pt, arc=2mm,
    left=2mm, right=2mm, top=1.5mm, bottom=1.5mm]
    \textbf{Primary Objective:} Students can describe a vector by its components and as an
    arrow in the plane, add and scale vectors (predicting the effect on the arrow), and
    compute a vector's length $\|\vv\|=\sqrt{v_1^{\,2}+v_2^{\,2}}$, explaining that length
    measures straight-line distance and is never negative.
\end{tcolorbox}

\renewcommand{\arraystretch}{1.4}
\begin{skillbox}[Priority Ideas \& Skills]{goldbox}
\begin{tabularx}{\linewidth}{@{}X|X@{}}
\textbf{Priority Ideas \& Skills} & \textbf{Key Understandings (the ``why'')} \\[2pt]
\hline
\vspace{-6pt}
\begin{itemize}[leftmargin=1.1em, nosep, after=\vspace{2pt}]
  \item Read/write a 2-D column vector and its components.
  \item Add vectors and multiply by a scalar (component rules).
  \item Compute magnitude with the Pythagorean theorem.
  \item Recognize a unit vector (length $1$).
\end{itemize}
&
\vspace{-6pt}
\begin{itemize}[leftmargin=1.1em, nosep, after=\vspace{2pt}]
  \item A vector carries \emph{two} ideas at once: direction and size.
  \item Adding is ``tip to tail''; scaling stretches/flips the same arrow.
  \item Length is distance, so it comes from Pythagoras and is $\ge 0$.
  \item These operations are the vocabulary the whole unit is built on.
\end{itemize}
\end{tabularx}
\end{skillbox}

\begin{skillbox}[{Vocabulary, Concepts, \& Theorems}]{greenbox}
\renewcommand{\arraystretch}{1.3}
\begin{tabularx}{\linewidth}{@{}l X@{}}
\textbf{Vector} & An ordered list of numbers; in 2-D, two components pictured as an arrow (or point). \\
\textbf{Scalar} & A single number --- size, no direction --- used to scale a vector. \\
\textbf{Components} & The entries $v_1,v_2$ of a vector; its coordinates. \\
\textbf{Vector addition} & \TallMath{\vv+\ww=\begin{bsmallmatrix} v_1+w_1 \\ v_2+w_2 \end{bsmallmatrix}}; geometrically tip-to-tail. \\
\textbf{Scalar multiplication} & \TallMath{c\,\vv=\begin{bsmallmatrix} c\,v_1 \\ c\,v_2 \end{bsmallmatrix}}; stretch, shrink, or flip. \\
\textbf{Magnitude (length)} & \TallMath{\|\vv\|=\sqrt{v_1^{\,2}+v_2^{\,2}}}; straight-line length, never negative. \\
\textbf{Unit vector} & A vector of length $1$. \\
\textbf{Linear combination} & $c\,\vv + d\,\ww$ --- scale then add (previewed; the focus of 1.1). \\
\end{tabularx}
\end{skillbox}

\begin{fixedskillbox}[Activate Prior Knowledge \& Spiral Review]{blush}
    The Warm-Up rehearses the three prerequisite skills this lesson leans on:
    \textbf{(1)} signed-number arithmetic (for component addition and scaling),
    \textbf{(2)} the Pythagorean theorem (which \emph{is} the magnitude formula), and
    \textbf{(3)} reading an ordered pair off the plane (which becomes reading a vector's
    components). Debrief item~2 out loud --- students should see the hypotenuse computation
    return as $\|\vv\|$ within the first ten minutes.
\end{fixedskillbox}

\begin{skillbox}[Hook]{blush}
    Project the drone flying $3$ east and $4$ north, recorded as $(3,4)$. Ask: \textbf{``One
    pair of numbers stored the entire trip --- what single number tells us the straight-line
    distance home, and how would you get it from the $3$ and the $4$?''} Let a few students
    commit to ``$5$, using the triangle'' before naming it magnitude. This motivates both the
    vector (a trip in two numbers) and its length (Pythagoras) in one picture.
\end{skillbox}

\begin{skillbox}[Lesson]{blush}
\begin{multicols}{2}
\textbf{1. What a vector is.} Define a vector as an ordered list; in 2-D, two
\textbf{components}. Show both pictures --- point $(v_1,v_2)$ and arrow from the origin ---
using $\vv=\begin{bsmallmatrix} 3\\4 \end{bsmallmatrix}$. Contrast with a \textbf{scalar}
(a lone number). Have students say aloud ``$3$ east, $4$ north.''

\textbf{2. Adding and scaling.} Give the component rules, then tie each to a picture:
addition is \textbf{tip-to-tail}; scaling by $c$ \textbf{stretches} ($c>1$),
\textbf{shrinks} ($0<c<1$), or \textbf{flips} ($c<0$) the same arrow. Work
$\vv+\ww$ and $2\vv,\ -\vv$ with $\ww=\begin{bsmallmatrix} 1\\-2 \end{bsmallmatrix}$. Name
$c\vv+d\ww$ a \textbf{linear combination} and flag it as the unit's throughline.

\columnbreak

\textbf{3. Length (magnitude).} Draw the right triangle under the arrow and derive
$\|\vv\|=\sqrt{v_1^2+v_2^2}$ from Pythagoras. Compute
$\left\|\begin{bsmallmatrix} 3\\4 \end{bsmallmatrix}\right\|=5$ and connect back to the Hook.
Stress two ideas: length is a \textbf{scalar}, and it is \textbf{never negative} (the squares
erase the signs). Introduce the \textbf{unit vector} (length $1$) and dividing by $\|\vv\|$.

\textbf{4. Guided practice.} Do the notes' practice set ($\vv=\begin{bsmallmatrix} 6\\8 \end{bsmallmatrix}$,
$\ww=\begin{bsmallmatrix} -3\\0 \end{bsmallmatrix}$) together, narrating the compute
$\rightarrow$ interpret loop: ``what did we get, and what does it mean?''
\end{multicols}
\end{skillbox}

\begin{skillbox}[Active Monitoring]{redbox}
Circulate during the notes practice and Tier work. Watch for and correct:
\begin{itemize}[leftmargin=1.2em, nosep]
  \item adding a scalar to a vector, or adding across components ($v_1+v_2$) instead of matching them;
  \item sign slips in scaling (e.g.\ $-\vv$) and in components like $\begin{bsmallmatrix} -1\\2 \end{bsmallmatrix}$;
  \item reporting a \emph{negative} length, or forgetting the square root.
\end{itemize}
Cold-call prompts: ``Point to $\vv+\ww$ on the grid.'' ``Why can't a length come out negative?''
``What does $\|\vv\|$ measure in the drone story?''
\end{skillbox}

\begin{skillbox}[Group Work \& Differentiation]{redbox}
\textbf{Drone Delivery} activity (tiered; mirrors the notes).
\begin{multicols}{3}
\textbf{Tier R --- Remediate.} Two given hops: describe a hop in words, add them, double one,
find one length. Rebuilds each definition on concrete integers.

\columnbreak
\textbf{Tier A --- Approaching.} Three-hop delivery: sum the hops, read off east/north, find
straight-line distance, and \emph{interpret} why distance travelled $\ne$ displacement.

\columnbreak
\textbf{Tier E --- Extension.} Build a unit vector in a given direction; test and justify when
$\|\vv+\ww\|=\|\vv\|+\|\ww\|$ (triangle inequality, informally).
\end{multicols}
\end{skillbox}

\begin{skillbox}[Individual Work \& Assessment]{redbox}
\textbf{Exit Ticket} (independent, no notes): add and scale a vector; compute a magnitude
($\|(5,12)\|=13$); and a \textbf{conceptual/justification check} --- explain why
$\begin{bsmallmatrix} -7\\0 \end{bsmallmatrix}$ does not have length $-7$ and give the correct
value ($7$). Collect and sort into ``got it / component slip / length misconception'' to decide
whether 1.1 opens with a quick re-teach.
\end{skillbox}

\begin{skillbox}[Reinforcement \& Extension]{goldbox}
\textbf{Homework:} mixed practice on adding/scaling and computing lengths (including simplest
radical form and a unit-vector check), plus a hiker displacement problem with an interpret
part. \textbf{Extension:} a vector as a \emph{data record} (a smoothie recipe scaled by a
batch factor) --- previews Unit~1's data theme. \textbf{Preview:} Lesson~1.1, \emph{Linear
Combinations of Vectors} --- doing both operations at once, $c\,\vv+d\,\ww$, and which points
in the plane two vectors can reach.
\end{skillbox}
\end{document}
